% GENERATED FILE. Edit canonical lessons, then run npm run generate:books.
% canonical-source-hash: fnv1a64:37d576d72099f1e4
% canonical-lessons: HI-C517-it, HI-C517-cuna, HI-C517-palastar, HI-C517-lipna, HI-C517-cinna

\chapter{Bricks and Mortar}
\label{ch:hi-bricks-and-mortar}
\hlchaptermodality{517}

\begin{chapteropening}
\textbf{By the end of this chapter:} \emph{I can talk about building and mending walls.}

Complete the last lesson of chapter 517: I can talk about building and mending walls.
\end{chapteropening}

\section[\texorpdfstring{\hi{ईंट} (ī̃ṭ)}{ईंट (ī̃ṭ)}]{\hi{ईंट} (ī̃ṭ) — a brick}

\label{lesson:HI-C517-it}



\begin{quote}
Before the new one: say the Hindi for a hole, then the Hindi for to trickle down.
\end{quote}

\subsection*{You'll want to know: \hi{ईंट}}
\textbf{\hi{ईंट}} — \emph{ī̃ṭ} — \textquotedblleft{}a brick\textquotedblright{}.

\begin{quote}
\hi{मज़दूर} \hi{ईंटें} \hi{सिर} \hi{पर} \hi{उठाकर} \hi{ले} \hi{जा} \hi{रहे} \hi{हैं।}
\end{quote}

\begin{quote}
\emph{Mazdūr ī̃ṭeṁ sir par uṭhākar le jā rahe haiṁ.}
\end{quote}

\begin{quote}
The labourers are carrying bricks on their heads.
\end{quote}

Feminine.

\begin{cousinweb}
\textbf{\hi{ईंट}} (\emph{ī̃ṭ}) is a tadbhava, inherited from Sanskrit \emph{iṣṭakā}, \textquotedblleft{}a brick\textquotedblright{} (CDIAL 1600).
\end{cousinweb}

\begin{grammarlens}[title={What you've built}]
One more word for describing what you lived and felt.
\end{grammarlens}

\subsection*{Your turn}
\begin{itemize}
\raggedright
  \item \emph{Say it:} \emph{ī̃ṭ}
  \item \emph{Say it:} \emph{ī̃ṭ}, once more
  \item \emph{Say it:} \emph{Mazdūr ī̃ṭeṁ sir par uṭhākar le jā rahe haiṁ.} — The labourers are carrying bricks on their heads.
  \item \emph{Recall:} say the Hindi for a hole, then the Hindi for to trickle down, then say \emph{ī̃ṭ} again
\end{itemize}

\begin{tcolorbox}[breakable,colback=teal!4,colframe=teal!35!black,title={Before you move on}]
What does \hi{ईंट} mean? (\textquotedblleft{}A brick\textquotedblright{}.) Say it once more.
\end{tcolorbox}
\section[\texorpdfstring{\hi{चूना} (cūnā)}{चूना (cūnā)}]{\hi{चूना} (cūnā) — lime (for whitewash)}

\label{lesson:HI-C517-cuna}



\begin{quote}
Before the new one: say the Hindi for to trickle down, then the Hindi for a brick.
\end{quote}

\subsection*{You'll want to know: \hi{चूना}}
\textbf{\hi{चूना}} — \emph{cūnā} — \textquotedblleft{}lime (for whitewash)\textquotedblright{}.

\begin{quote}
\hi{दीवारों} \hi{पर} \hi{नया} \hi{चूना} \hi{लगा} \hi{है}, \hi{इसलिए} \hi{कमरा} \hi{साफ़} \hi{दिखता} \hi{है।}
\end{quote}

\begin{quote}
\emph{Dīvāroṁ par nayā cūnā lagā hai, isliye kamrā sāf dikhtā hai.}
\end{quote}

\begin{quote}
The walls have fresh whitewash, so the room looks clean.
\end{quote}

Masculine.

\begin{cousinweb}
\textbf{\hi{चूना}} (\emph{cūnā}) is a tadbhava, inherited from Sanskrit \emph{cūrṇa}, \textquotedblleft{}powder\textquotedblright{}, which already meant lime in Sanskrit (CDIAL 4889).
\end{cousinweb}

\begin{grammarlens}[title={What you've built}]
One more word for describing what you lived and felt.
\end{grammarlens}

\subsection*{Your turn}
\begin{itemize}
\raggedright
  \item \emph{Say it:} \emph{cūnā}
  \item \emph{Say it:} \emph{cūnā}, once more
  \item \emph{Say it:} \emph{Dīvāroṁ par nayā cūnā lagā hai, isliye kamrā sāf dikhtā hai.} — The walls have fresh whitewash, so the room looks clean.
  \item \emph{Recall:} say the Hindi for to trickle down, then the Hindi for a brick, then say \emph{cūnā} again
\end{itemize}

\begin{tcolorbox}[breakable,colback=teal!4,colframe=teal!35!black,title={Before you move on}]
What does \hi{चूना} mean? (\textquotedblleft{}Lime (for whitewash)\textquotedblright{}.) Say it once more.
\end{tcolorbox}
\section[\texorpdfstring{\hi{पलस्तर} (palastar)}{पलस्तर (palastar)}]{\hi{पलस्तर} (palastar) — plaster (on a wall)}

\label{lesson:HI-C517-palastar}



\begin{quote}
Before the new one: say the Hindi for a brick, then the Hindi for lime (for whitewash).
\end{quote}

\subsection*{You'll want to know: \hi{पलस्तर}}
\textbf{\hi{पलस्तर}} — \emph{palastar} — \textquotedblleft{}plaster (on a wall)\textquotedblright{}.

\begin{quote}
\hi{दीवार} \hi{पर} \hi{नया} \hi{पलस्तर} \hi{अभी} \hi{गीला} \hi{है।}
\end{quote}

\begin{quote}
\emph{Dīvār par nayā palastar abhī gīlā hai.}
\end{quote}

\begin{quote}
The new plaster on the wall is still wet.
\end{quote}

Masculine.

\begin{cousinweb}
\textbf{\hi{पलस्तर}} (\emph{palastar}) is a loan from English \emph{plaster}.
\end{cousinweb}

\begin{grammarlens}[title={What you've built}]
One more word for describing what you lived and felt.
\end{grammarlens}

\subsection*{Your turn}
\begin{itemize}
\raggedright
  \item \emph{Say it:} \emph{palastar}
  \item \emph{Say it:} \emph{palastar}, once more
  \item \emph{Say it:} \emph{Dīvār par nayā palastar abhī gīlā hai.} — The new plaster on the wall is still wet.
  \item \emph{Recall:} say the Hindi for a brick, then the Hindi for lime (for whitewash), then say \emph{palastar} again
\end{itemize}

\begin{tcolorbox}[breakable,colback=teal!4,colframe=teal!35!black,title={Before you move on}]
What does \hi{पलस्तर} mean? (\textquotedblleft{}Plaster (on a wall)\textquotedblright{}.) Say it once more.
\end{tcolorbox}
\section[\texorpdfstring{\hi{लीपना} (līpnā)}{लीपना (līpnā)}]{\hi{लीपना} (līpnā) — to smear (a floor or wall with mud)}

\label{lesson:HI-C517-lipna}



\begin{quote}
Before the new one: say the Hindi for lime (for whitewash), then the Hindi for plaster (on a wall).
\end{quote}

\subsection*{You'll want to know: \hi{लीपना}}
\textbf{\hi{लीपना}} — \emph{līpnā} — \textquotedblleft{}to smear (a floor or wall with mud)\textquotedblright{}.

\begin{quote}
\hi{गाँव} \hi{में} \hi{औरतें} \hi{मिट्टी} \hi{से} \hi{आँगन} \hi{लीपती} \hi{हैं।}
\end{quote}

\begin{quote}
\emph{Gā̃v meṁ aurteṁ miṭṭī se āṁgan līptī haiṁ.}
\end{quote}

\begin{quote}
In the village the women smear the courtyard with mud.
\end{quote}

Past, with \hi{ने} (he or she …): \hi{उसने} \hi{चूल्हा} \hi{लीपा}, \emph{usne cūlhā līpā}.

\begin{cousinweb}
\textbf{\hi{लीपना}} (\emph{līpnā}) is a tadbhava, inherited from Sanskrit \emph{lipyatē}, \textquotedblleft{}is smeared\textquotedblright{} (CDIAL 11061).
\end{cousinweb}

\begin{grammarlens}[title={What you've built}]
One more word for describing what you lived and felt.
\end{grammarlens}

\subsection*{Your turn}
\begin{itemize}
\raggedright
  \item \emph{Say it:} \emph{līpnā}
  \item \emph{Say it:} \emph{līpnā}, once more
  \item \emph{Say it:} \emph{Gā̃v meṁ aurteṁ miṭṭī se āṁgan līptī haiṁ.} — In the village the women smear the courtyard with mud.
  \item \emph{Recall:} say the Hindi for lime (for whitewash), then the Hindi for plaster (on a wall), then say \emph{līpnā} again
\end{itemize}

\begin{tcolorbox}[breakable,colback=teal!4,colframe=teal!35!black,title={Before you move on}]
What does \hi{लीपना} mean? (\textquotedblleft{}To smear (a floor or wall with mud)\textquotedblright{}.) Say it once more.
\end{tcolorbox}
\section[\texorpdfstring{\hi{चिनना} (cinnā)}{चिनना (cinnā)}]{\hi{चिनना} (cinnā) — to lay (bricks); to build up (a wall)}

\label{lesson:HI-C517-cinna}



\begin{quote}
Before the new one: say the Hindi for plaster (on a wall), then the Hindi for to smear (a floor or wall with mud).
\end{quote}

\subsection*{You'll want to know: \hi{चिनना}}
\textbf{\hi{चिनना}} — \emph{cinnā} — \textquotedblleft{}to lay (bricks); to build up (a wall)\textquotedblright{}.

\begin{quote}
\hi{मिस्त्री} \hi{ईंटें} \hi{चिनकर} \hi{दीवार} \hi{बना} \hi{रहा} \hi{है।}
\end{quote}

\begin{quote}
\emph{Mistrī ī̃ṭeṁ cinkar dīvār banā rahā hai.}
\end{quote}

\begin{quote}
The mason is laying bricks and building a wall.
\end{quote}

Past, with \hi{ने} (he or she …): \hi{उसने} \hi{दीवार} \hi{चिनी}, \emph{usne dīvār cinī}.

\begin{cousinweb}
\textbf{\hi{चिनना}} (\emph{cinnā}) is a tadbhava, inherited from Sanskrit \emph{cinōti}, \textquotedblleft{}heaps up, collects\textquotedblright{} (CDIAL 4814). CDIAL gives the Hindi verb as \textquotedblleft{}to build up (a wall), gather\textquotedblright{}.
\end{cousinweb}

\begin{grammarlens}[title={What you've built}]
One more word for describing what you lived and felt.
\end{grammarlens}

\subsection*{Your turn}
\begin{itemize}
\raggedright
  \item \emph{Say it:} \emph{cinnā}
  \item \emph{Say it:} \emph{cinnā}, once more
  \item \emph{Say it:} \emph{Mistrī ī̃ṭeṁ cinkar dīvār banā rahā hai.} — The mason is laying bricks and building a wall.
  \item \emph{Recall:} say the Hindi for plaster (on a wall), then the Hindi for to smear (a floor or wall with mud), then say \emph{cinnā} again
\end{itemize}

\begin{tcolorbox}[breakable,colback=teal!4,colframe=teal!35!black,title={Before you move on}]
What does \hi{चिनना} mean? (\textquotedblleft{}To lay (bricks); to build up (a wall)\textquotedblright{}.) Say it once more.
\end{tcolorbox}
